\originalchapter{用光和轨道测量恒星}
8.901没有公开完整课堂讲义。本章及后续恒星专题按官方26次课历和完整公开作业独立组织，不能称作逐堂课堂笔录。本章支撑第1--4次课：天体不能直接称重，距离、光谱、双星运动如何共同给出物理尺度，是学习后面恒星结构方程的入口。需要Newton力学、圆锥曲线、对数和角度的小量近似。

\originalsection{天空坐标先描述方向，距离另需标尺}
赤经$\alpha$、赤纬$\delta$在赤道坐标系指定一个方向；赤经常用时、分、秒且24小时对应$2\pi$，赤纬用角度。银经、银纬则相对银河平面定义，转换是两套坐标轴之间的旋转，不是改变物体距离。自转使本地地平坐标随时刻变化，地球轨道、历元与自行又影响精密位置；仅靠一个天空角度不能给三维空间位置。

年视差$p$来自观测基线1 AU：小角下$d={\rm AU}/p_{\rm rad}$。定义秒差距使$d/{\rm pc}=1/(p/{\rm arcsec})$。这里的$p$是视差椭圆的半轴尺度，不能将半年最大角位移$2p$再当成$p$。对小误差$\sigma_d/d\simeq\sigma_p/p$；误差大或视差非正时，直接倒数会产生非对称偏差，必须用测量似然与距离先验，不把负视差解释为负距离。

光度和波段星等的定义见第\ref{ch:radiation}章。无消光的低红移距离模数为$m-M=5\log_{10}(d/10{\rm pc})$；有消光则多$A_\lambda$。观测通量只给$L/d^2$，若不知道距离或绝对光度，二者互相简并。光谱型约束大气温度，光度级约束表面重力与线宽，不能单由一种颜色同时唯一解距离、温度、半径和尘埃。

\originalsection{H--R图把两个可测维度连接到结构}
H--R图的纵轴为光度或绝对星等，横轴为有效温度或谱型；温度惯例向右降低。有效温度通过$L=4\pi R^2\sigma_{\rm SB}T_{\rm eff}^4$定义，即同光度、同面积黑体的温度，不等于真实谱在每个波段都黑体。给$L,T_{\rm eff}$即可求$R$；同温度的巨星比主序星亮，必须有更大半径。

主序是一族正在核心烧氢的平衡模型，不是某颗星永远沿同一条线移动。白矮星的高温、低光度说明面积小；红巨星低表面温度仍可高光度，因为半径很大。色--星等图还受滤光片、距离、消光与总光度校正影响，只有在这些量受控时才可当成H--R图的直接代理。星团星基本同距、近同龄，有助于把离开主序的质量连到年龄；尚需金属丰度、双星和测量选择的模型。

\originalsection{两体问题为何能称重}
分离质心后，相对向量$\boldsymbol r=\boldsymbol r_2-\boldsymbol r_1$满足$\ddot{\boldsymbol r}=-G(m_1+m_2)\boldsymbol r/r^3$。单位约化质量的角动量$\ell=r^2\dot\phi$守恒。令$u=1/r$，利用$\dd/\dd t=\ell u^2\dd/\dd\phi$，径向方程化成
\[
 u''+u=\frac{G(m_1+m_2)}{\ell^2}.
\]
选择近心方向为$\phi=0$，解为$r=p/(1+e\cos\phi)$，$p=\ell^2/[G(m_1+m_2)]$。能量符号决定椭圆、抛物线或双曲线；$e<1$时$p=a(1-e^2)$。

面积速度$\dot A=\ell/2$，椭圆面积$\pi a^2\sqrt{1-e^2}$。周期等于面积除面积速度：
\[
 P=\frac{2\pi a^2\sqrt{1-e^2}}{\ell},\qquad
 P^2=\frac{4\pi^2a^3}{G(m_1+m_2)}.
\]
所以观测相对轨道的\emph{实际}半长轴与周期可给总质量。目视双星只给角半长轴$\alpha_a$，还需距离使$a=d\alpha_a$。绕质心两半长轴之比$a_1/a_2=m_2/m_1$又给质量比。质量与距离因此互相进入；不能拿投影瞬时间距直接替代真实轨道半长轴。

\originalsection{光谱双星的质量函数保留了倾角}
视向速度是沿观测方向的投影。椭圆轨道拟合给周期、偏心率、速度半振幅$K_1$，有
\[
 K_1=\frac{2\pi a_1\sin i}{P\sqrt{1-e^2}},\qquad
 f(m)=\frac{PK_1^3(1-e^2)^{3/2}}{2\pi G}
 =\frac{m_2^3\sin^3i}{(m_1+m_2)^2}.
\]
此式由$a_1=a\,m_2/(m_1+m_2)$和Kepler律消掉$a$，不是另设经验质量关系。单线光谱看不到伴星速度，$f(m)$不是伴星质量；$\sin i\le1$只给依赖$m_1$的下限。双线光谱可得$m_2/m_1=K_1/K_2$，若还有食和轨道几何就能约束$i$，解除更多简并。8.901 PS1.9的质量函数分布题会实际展示为何极端值和随机倾角检验不能忽略样本选择。

食双星的接触时刻和光变形状还给相对半径$(R_1+R_2)/a$、表面亮度比及倾角。掩食深度并不总等于半径平方比，因为两星温度不同、边缘昏暗、部分食和光污染都会改变它。轨道测质量，大气光谱测温度，距离和总通量测半径，是独立信息相互校验的物理链。
